\documentclass[a4paper,12pt]{article}

% --------------------------------------------------------
% CODIFICAÇÃO E TIPOGRAFIA
% --------------------------------------------------------
\usepackage[utf8]{inputenc}       % Permite acentuação direta no código-fonte
\usepackage[T1]{fontenc}          % Usa codificação T1 (melhor para português)
\usepackage[brazil]{babel}        % Tradução e hifenização para português
\usepackage{microtype}            % Melhora o espaçamento entre letras/palavras
\usepackage{parskip}              % Espaço entre parágrafos em vez de recuo

% --------------------------------------------------------
% PACOTES MATEMÁTICOS
% --------------------------------------------------------
\usepackage{amsmath}              % Ambientes matemáticos como align, equation
%\usepackage{amssymb}              % Símbolos matemáticos adicionais
\usepackage{amsthm}               % Ambientes de teoremas
\usepackage{mathtools}            % Extensões do amsmath (ex: \coloneqq)
\usepackage{mathrsfs}             % Fonte caligráfica com \mathscr
\usepackage{bbm}                  % Indicador \mathbbm{1}
\usepackage{dsfont}               % Alternativa para conjuntos: \mathds
\usepackage{commath}              % Notações como \abs{}, \norm{}, \eval{}
\usepackage{pgfmath}             % Permite realizar cálculos matemáticos (somas, produtos, números aleatórios, etc.)
\usepackage{siunitx}              % sistema de unidades
\sisetup{group-separator = {\,}} % espaço fino

% --------------------------------------------------------
% FONTE MODERNA (TEXTO E MATEMÁTICA)
% --------------------------------------------------------
%\usepackage{newtxtext}            % Fonte do texto (estilo Times)
%\usepackage{newtxmath}            % Fonte matemática compatível com pacotes AMS
\usepackage[bitstream-charter]{mathdesign} % Fonte elegante e matemática harmoniosa

% --------------------------------------------------------
% AMBIENTES TEÓRICOS PERSONALIZADOS
% --------------------------------------------------------
\theoremstyle{definition}
\newtheorem{definicao}{Definição}       % Definição numerada por seção
\newtheorem{exemplo}{Exemplo}           % Exemplo numerado por seção

%\theoremstyle{plain}
\newtheorem{teorema}{Teorema}           % Teorema numerado por seção
\newtheorem{lema}[teorema]{Lema}                 % Lema com numeração conjunta
\newtheorem{proposicao}[teorema]{Proposição}     % Proposição idem
\newtheorem{corolario}[teorema]{Corolário}       % Corolário idem
\newtheorem{exercicio}[teorema]{Exercício}

%\theoremstyle{remark}
\newtheorem{observacao}[teorema]{Observação}     % Observação com mesmo contador

% --------------------------------------------------------
% TABELAS E FIGURAS
% --------------------------------------------------------
\usepackage{graphicx}             % Inclusão de imagens
\usepackage{float}                % Controle de posicionamento (ex: [H])
\usepackage{caption}              % Personalização de legendas
\usepackage{subcaption}          % Subfiguras com \begin{subfigure}
\usepackage{booktabs}            % Tabelas com qualidade profissional
\usepackage{array}               % Mais opções em colunas de tabelas
\usepackage{multirow}            % Células que ocupam várias linhas
\usepackage{multicol}            % Disposição em colunas múltiplas
\usepackage{tabularx}            % Tabelas com colunas de largura ajustável
\usepackage{longtable}           % Tabelas que quebram página
\usepackage{arydshln}            % Linhas tracejadas horizontais e verticais
\usepackage{makecell}            % Células com múltiplas linhas (e quebra de linha)
\usepackage{diagbox}             % Cabeçalho diagonal em tabelas
\usepackage{pdflscape}           % Gira páginas (modo paisagem)
\newcolumntype{L}[1]{>{\raggedright\let\newline\\\arraybackslash\hspace{0pt}}m{#1}}
\newcolumntype{C}[1]{>{\centering\let\newline\\\arraybackslash\hspace{0pt}}m{#1}}
\newcolumntype{R}[1]{>{\raggedleft\let\newline\\\arraybackslash\hspace{0pt}}m{#1}}

% --------------------------------------------------------
% FORMATAÇÃO E ESTRUTURA
% --------------------------------------------------------
\usepackage{geometry}            % Configuração de margens
\geometry{a4paper, left=1cm, right=1cm, bottom=1.5cm, top=1.5cm, headsep=1cm, footskip=1cm}
\usepackage{titlesec}            % Controle sobre títulos de seções
\usepackage{fancyhdr}            % Cabeçalhos e rodapés personalizados
\pagestyle{fancy}
\fancyhf{}                % Limpa cabeçalho e rodapé
\fancyfoot[R]{\thepage}   % Numeração à direita no rodapé
\renewcommand{\headrulewidth}{0pt} % Remove linha no topo
\usepackage[shortlabels]{enumitem}            % Listas com controle de espaçamento e símbolo
\usepackage{tikz}                % Desenho de gráficos vetoriais (diagramas, etc.)
\usetikzlibrary{matrix,arrows,patterns,snakes,decorations.pathreplacing,3d,arrows.meta,calc}
\usepackage{etoolbox}             % Permite lógica condicional e manipulação de comandos
\usepackage{background}           % Permite adicionar conteúdo fixo no fundo da página (ex: linhas, marcas d'água)
\SetBgContents{} % Remove conteúdo padrão ("DRAFT")
\definecolor{background-color}{RGB}{233 227 206}

% --------------------------------------------------------
% GRÁFICOS E PLOTAGENS
% --------------------------------------------------------
\usepackage{pgfplots}             % Criação de gráficos vetoriais em 2D e 3D
\pgfplotsset{compat=1.18}         % Compatibilidade com versão atual do pacote
\usepackage{currfile} % Permite saber qual é o caminho completo do arquivo atual,

% --------------------------------------------------------
% LINKS E CORES
% --------------------------------------------------------
\usepackage{xcolor}              % Definição de cores (texto, links, etc.)
\usepackage[hidelinks]{hyperref} % Links sem moldura colorida no PDF
\usepackage{url}                 % Quebra automática de URLs
% Dica: se quiser links coloridos, use:
% \usepackage[colorlinks=true, linkcolor=blue, citecolor=blue, urlcolor=blue]{hyperref}

% --------------------------------------------------------
% CAIXAS E DESTAQUES VISUAIS
% --------------------------------------------------------
\usepackage[most]{tcolorbox}      % Criação de caixas coloridas e personalizadas para destaques, avisos, exemplos, etc.

% --------------------------------------------------------
% CÓDIGOS E ALGORITMOS
% --------------------------------------------------------
\usepackage{listings}            % Inclusão de código-fonte com destaque
\usepackage{algorithm}           % Ambiente de algoritmo
\usepackage[noend]{algpseudocode} % Sintaxe tipo pseudocódigo

% --------------------------------------------------------
% COMANDOS
% --------------------------------------------------------

\newcommand{\bbN}{\mathbb{N}}
\newcommand{\bbZ}{\mathbb{Z}}
\newcommand{\bbQ}{\mathbb{Q}}
\newcommand{\bbR}{\mathbb{R}}
\newcommand{\calnN}{\mathcal{N}}

\newcommand{\assunto}{FUNÇÕES DE DISTRIBUIÇÃO}
\newcommand{\atividade}{lista} % prova ou lista
\newcommand{\turma}{ESTATÍSTICA}
\newcommand{\professor}{Igor Martins Silva}
\newcommand{\semestre}{2}
\newcommand{\ano}{2026}
\newcommand{\mesTexto}{outubro}
\newcommand{\mesNum}{10}
\newcommand{\dia}{14}
\newcommand{\dataTexto}{\dia\ de \mesTexto\ de \ano}
\newcommand{\dataNun}{\dia/\mesNum/\ano}

\ifdefstring{\atividade}{prova}{
    \newcommand{\titulo}{Avaliação de Estatística}
}{
    \newcommand{\titulo}{Lista de exercícios de Estatística}
}

\edef\pathCapa{\currfiledir}

\newcommand{\subtitulo}{
    %
    \noindent
    \begin{minipage}{2.8cm}
        \includegraphics[width=2.8cm]{\pathCapa Logo_Cefet.png}
    \end{minipage}
    \hfill
    \begin{minipage}{\dimexpr\textwidth-6.2cm\relax}
        \begin{center}
            CEFET-MG -- Campus Contagem \\[3pt]
            \titulo \ -- \semestre.\textordmasculine\ Semestre de \ano \\[3pt]
            \professor
        \end{center}
    \end{minipage}
    \hfill
    \begin{minipage}{3.2cm}
        \includegraphics[width=3.2cm]{\pathCapa Brasao_Brasil.png}
    \end{minipage}
    %
    
    %
    \begin{center}
        \vspace{15pt}
        \begin{tabular}{|C{250pt}|C{250pt}|}
            \hline
            \textbf{ASSUNTO} &
            \textbf{TURMA} \\
            \hline
            \assunto & \turma \\
            \hline
        \end{tabular}
    \end{center}
    %
    \vspace{15pt}
}

\newcommand{\subtituloAlt}{
    %
    \noindent
    \begin{minipage}{2.8cm}
        \includegraphics[width=2.8cm]{\pathCapa Logo_Cefet.png}
    \end{minipage}
    \hfill
    \begin{minipage}{\dimexpr\textwidth-6.2cm\relax}
        \begin{center}
            CEFET-MG -- Campus Contagem \\[3pt]
            \titulo \ -- \semestre.\textordmasculine\ Semestre de \ano \\[3pt]
            \professor
        \end{center}
    \end{minipage}
    \hfill
    \begin{minipage}{3.2cm}
        \includegraphics[width=3.2cm]{\pathCapa Brasao_Brasil.png}
    \end{minipage}
    %
    
    %
    \begin{center}
        \vspace{15pt}
        \begin{tabular}{|C{250pt}|C{100pt}|C{100pt}|}
            \hline
            \textbf{ASSUNTO} &
            \textbf{DATA} &
            \textbf{VALOR} \\
            \hline
            \assunto & \dataNun & \ValorProva\ pontos \\
            \hline
        \end{tabular}
    \end{center}
    %
    \vspace{15pt}
}

\newcommand{\valorquestao}{}

\newcommand{\definevalorquestao}{%
    \ifdefstring{\atividade}{lista}{%
        % Não faz nada para lista
    }{%
        \renewcommand{\valorquestao}{\ValorQ pontos}
    }%
}

\newenvironment{exercicioBanco}[1][]{%
    \ifdefstring{\atividade}{prova}{%
        \begin{exercicio}[#1]%
    }{%
        \begin{exercicio}%
    }%
}{%
    \end{exercicio}
}

\ifdefstring{\atividade}{lista}{
    \pagecolor{background-color}
}{
    
}



\hypersetup{
    pdfauthor={\professor},
    pdftitle={\titulo -- \assunto},
}

\title{\titulo -- \assunto}
\author{\professor}
\date{\data}

%************* INÍCIO DO DOCUMENTO *************
\begin{document}
    \subtitulo
        
    \phantomsection
    \addcontentsline{toc}{section}{Exercício 1}
    %%%%%%%%%%%%%%%%%%%%%%%%%%%%%%%%%
%
%       EXERCÍCIO
%
%%%%%%%%%%%%%%%%%%%%%%%%%%%%%%%%%

\definevalorquestao

\begin{exercicioBanco}[\valorquestao]
Uma variável aleatória discreta \(X\) tem a seguinte função de distribuição
%
\[
    F(x) = \left\{
    \begin{array}{rl}
        0, & \text{se } x < -1, \\
        0,2, & \text{se } -1 \le x < 2, \\
        0,5, & \text{se } 2 \le x < 5, \\
        0,7, & \text{se } 5 \le x < 6, \\
        0,9, & \text{se } 6 \le x < 15, \\
        1 & \text{se } x \ge 15. \\
    \end{array}
    \right.
\]
%
Determine:
\begin{enumerate}[(a)]
    \item A função de probabilidade de \(X\).
    %
    \item \(P(X \le -2)\).
    %
    \item \(P(X < 2)\).
    %
    \item \(P(3 \le X \le 12)\).
\end{enumerate}
%
\end{exercicioBanco}


    \noindent\rule{\linewidth}{1pt}  % linha horizontal fina
    
    \phantomsection
    \addcontentsline{toc}{section}{Exercício 2}
    %%%%%%%%%%%%%%%%%%%%%%%%%%%%%%%%%
%
%       EXERCÍCIO
%
%%%%%%%%%%%%%%%%%%%%%%%%%%%%%%%%%

\definevalorquestao

\begin{exercicioBanco}[\valorquestao]
Em um certo restaurante, paga-se pelo almoço uma quantia fixa dependendo da escolha feita de prato e bebida. A carne de peixe tem 10\% de preferência, enquanto frango tem 40\% e carne bovina 50\%. As três escolhas bebidas estão condicionadas à opção do prato, segundo a tabela abaixo:
%
\[
    \begin{array}{c|c|c|c}
         & \text{Cerveja} & \text{Água} & \text{Vinho} \\
         \hline
         P(\text{Bebida}\mid\text{Peixe}) & 0,4 & 0,3 & 0,3 \\
         P(\text{Bebida}\mid\text{Frango}) & 0,3 & 0,5 & 0,2 \\
         P(\text{Bebida}\mid\text{Bovina}) & 0,6 & 0,3 & 0,1 \\
    \end{array}
\]
%
Admita os seguintes preços
%
\[
    \begin{array}{c|c|c|c|c|c|c}
        \text{Pedido} & \text{Peixe} & \text{Frango} & \text{Bovina} & \text{Cerveja} & \text{Água} & \text{Vinho} \\ 
        \hline
        \text{Preço} & 12 & 15 & 18 & 6 & 3 & 9
    \end{array}
\]
%
Determine a função de probabilidade para a variável \(X:\) preço do almoço.
\end{exercicioBanco}


    \noindent\rule{\linewidth}{1pt}  % linha horizontal fina
    
    \phantomsection
    \addcontentsline{toc}{section}{Exercício 3}
    %%%%%%%%%%%%%%%%%%%%%%%%%%%%%%%%%
%
%       EXERCÍCIO
%
%%%%%%%%%%%%%%%%%%%%%%%%%%%%%%%%%

\definevalorquestao

\begin{exercicioBanco}[\valorquestao]
A resistência (em toneladas) de vigas de concreto produzidas por uma empresa, comporta-se conforme a função de probabilidade abaixo:
%
\[
    \begin{array}{c|c|c|c|c|c}
         \text{Resistência} & 2 & 3 & 4 & 5 & 6 \\
         \hline
         p_{i} & 0,1 & 0,1 & 0,4 & 0,2 & 0,2
    \end{array}
\]
%
Admita que essas vigas são aprovadas para uso em construções se suportarem pelo menos \(3\) toneladas. De um grande lote fabricado pela empresa, escolhemos 15 vigas ao acaso. Qual é a probabilidade de:
%
\begin{enumerate}[(a)]
    \item Todas serem aptas para construções?
    \item No mínimo \(13\) serem aptas?
\end{enumerate}
%
\end{exercicioBanco}


    \noindent\rule{\linewidth}{1pt}  % linha horizontal fina
    
    \phantomsection
    \addcontentsline{toc}{section}{Exercício 4}
    %%%%%%%%%%%%%%%%%%%%%%%%%%%%%%%%%
%
%       EXERCÍCIO
%
%%%%%%%%%%%%%%%%%%%%%%%%%%%%%%%%%

\ifdefstring{\atividade}{prova}{%
    \renewcommand{\valorquestao}{\ValorQ\ pontos}
}%

\begin{exercicioBanco}[\valorquestao]
Suponha que uma variável aleatória contínua tenha densidade de probabilidade dada por
%
\[
    f(x) = \left\{
    \begin{array}{rl}
        \frac{1}{6}x + k, & \text{se } 0 < x < 3, \\
        0, & \text{caso contrário}.
    \end{array}
    \right.
\]
%
\begin{enumerate}[(a)]
    \item Qual é o valor de \(k\)?
    \item Quanto vale \(b\) tal que \(P(X > b) = \frac{5}{9}\)?
\end{enumerate}
%
\end{exercicioBanco}


    \noindent\rule{\linewidth}{1pt}  % linha horizontal fina
    
    \phantomsection
    \addcontentsline{toc}{section}{Exercício 5}
    %%%%%%%%%%%%%%%%%%%%%%%%%%%%%%%%%
%
%       EXERCÍCIO
%
%%%%%%%%%%%%%%%%%%%%%%%%%%%%%%%%%

\ifdefstring{\atividade}{prova}{%
    \renewcommand{\valorquestao}{\ValorQ\ pontos}
}%

\begin{exercicioBanco}[\valorquestao]
Um usuário de transporte coletivo chega pontualmente às 8 horas para pegar o seu ônibus. Devido ao trânsito caótico, a demora pode ser qualquer tempo entre 1 e 20 minutos (admita que o relógio ``pule'' de minuto em minuto).
%
\begin{enumerate}[(a)]
    \item Qual a probabilidade de demorar mais de 10 minutos?
    %
    \item Qual a probabilidade de demorar pelo menos 5, mas não mais de 10 minutos?
\end{enumerate}
%
\end{exercicioBanco}


    \noindent\rule{\linewidth}{1pt}  % linha horizontal fina
    
    \phantomsection
    \addcontentsline{toc}{section}{Exercício 6}
    %%%%%%%%%%%%%%%%%%%%%%%%%%%%%%%%%
%
%       EXERCÍCIO
%
%%%%%%%%%%%%%%%%%%%%%%%%%%%%%%%%%

\definevalorquestao

\begin{exercicioBanco}[\valorquestao]
A espessura de um flange em um componente de espaçonave é uma variável aleatória contínua uniformemente distribuída entre 0,95 e 1,05 milímetros. Qual o valor da espessura que é excedida por 90\% dos flanges?
\end{exercicioBanco}


    \noindent\rule{\linewidth}{1pt}  % linha horizontal fina
    
    \phantomsection
    \addcontentsline{toc}{section}{Exercício 7}
    %%%%%%%%%%%%%%%%%%%%%%%%%%%%%%%%%
%
%       EXERCÍCIO
%
%%%%%%%%%%%%%%%%%%%%%%%%%%%%%%%%%

\definevalorquestao

\begin{exercicioBanco}[\valorquestao]
Um time de futebol tem probabilidade 0,92 de vitória sempre que joga. Se o time atuar 4 vezes, determine a probabilidade de que vença exatamente 1 partida.
\end{exercicioBanco}


    \noindent\rule{\linewidth}{1pt}  % linha horizontal fina
    
    \phantomsection
    \addcontentsline{toc}{section}{Exercício 8}
    %%%%%%%%%%%%%%%%%%%%%%%%%%%%%%%%%
%
%       EXERCÍCIO
%
%%%%%%%%%%%%%%%%%%%%%%%%%%%%%%%%%

\definevalorquestao

\begin{exercicioBanco}[\valorquestao]
Uma fábrica produz peças eletrônicas, das quais 70\% estão dentro dos padrões de qualidade. Um lote com 20 peças é selecionado aleatoriamente para inspeção.
%
\begin{enumerate}[(a)]
    \item Qual é a probabilidade de exatamente 12 peças estarem dentro dos padrões?
    %
    \item Qual é a probabilidade de pelo menos 2 peças apresentarem defeito?
\end{enumerate}
%
\end{exercicioBanco}


    \noindent\rule{\linewidth}{1pt}  % linha horizontal fina
    
    \phantomsection
    \addcontentsline{toc}{section}{Exercício 9}
    %%%%%%%%%%%%%%%%%%%%%%%%%%%%%%%%%
%
%       EXERCÍCIO
%
%%%%%%%%%%%%%%%%%%%%%%%%%%%%%%%%%

\definevalorquestao

\begin{exercicioBanco}[\valorquestao]
Um jogador participa de um jogo em que, a cada rodada, tem probabilidade \(0{,}3\) de vencer, independentemente das rodadas anteriores. O jogo é repetido até que ele obtenha sua primeira vitória. Qual é a probabilidade de o jogador vencer pela primeira vez na 4.\textordfeminine\ rodada?
%
\end{exercicioBanco}


    \noindent\rule{\linewidth}{1pt}  % linha horizontal fina
    
    \phantomsection
    \addcontentsline{toc}{section}{Exercício 10}
    %%%%%%%%%%%%%%%%%%%%%%%%%%%%%%%%%
%
%       EXERCÍCIO
%
%%%%%%%%%%%%%%%%%%%%%%%%%%%%%%%%%

\definevalorquestao

\begin{exercicioBanco}[\valorquestao]
Uma moeda equilibrada é lançada sucessivamente, de modo independente, até que ocorra a primeira cara. Seja \(X\) a variável aleatória que conta o número de lançamentos anteriores à ocorrência de cara. Determine
%
\begin{enumerate}[(a)]
    \begin{multicols}{2}
        \item \(P(X \le 2)\).
        %
        \item \(P(X > 1)\).
    \end{multicols}
\end{enumerate}
%
\end{exercicioBanco}


    \noindent\rule{\linewidth}{1pt}  % linha horizontal fina
    
    \phantomsection
    \addcontentsline{toc}{section}{Exercício 11}
    %%%%%%%%%%%%%%%%%%%%%%%%%%%%%%%%%
%
%       EXERCÍCIO
%
%%%%%%%%%%%%%%%%%%%%%%%%%%%%%%%%%

\ifdefstring{\atividade}{prova}{%
    \renewcommand{\valorquestao}{\ValorQ\ pontos}
}%

\begin{exercicioBanco}[\valorquestao]
Um caminho para chegar a uma festa pode ser dividido em três etapas. Sem enganos, o trajeto é feito em 1 hora e 30 minutos. Se enganos acontecem na primeira etapa, acrescente 15 minutos ao tempo do trajeto. Para enganos na segunda etapa, o acréscimo é de 20 minutos e, para a terceira, 25 minutos. Admita que a probabilidade de engano é 0,15; 0,20 e 0,25 para a primeira, segunda e terceira etapas, respectivamente. É provável haver atraso na chegada à festa? Determine a probabilidade de haver atraso e o atraso não passar de 40 minutos.
\end{exercicioBanco}


    \noindent\rule{\linewidth}{1pt}  % linha horizontal fina
    
    \phantomsection
    \addcontentsline{toc}{section}{Exercício 12}
    %%%%%%%%%%%%%%%%%%%%%%%%%%%%%%%%%
%
%       EXERCÍCIO
%
%%%%%%%%%%%%%%%%%%%%%%%%%%%%%%%%%

\definevalorquestao

\begin{exercicioBanco}[\valorquestao]
Em momentos de pico, a chegada de aviões a um aeroporto se dá segundo o modelo de Poisson, com taxa (média) de 1 por minuto. Determine a probabilidade de 3 chegadas em um minuto qualquer do horário de pico.
\end{exercicioBanco}


    \noindent\rule{\linewidth}{1pt}  % linha horizontal fina
    
    \phantomsection
    \addcontentsline{toc}{section}{Exercício 13}
    %%%%%%%%%%%%%%%%%%%%%%%%%%%%%%%%%
%
%       EXERCÍCIO
%
%%%%%%%%%%%%%%%%%%%%%%%%%%%%%%%%%

\ifdefstring{\atividade}{prova}{%
    \renewcommand{\valorquestao}{\ValorQ\ pontos}
}%

\begin{exercicioBanco}[\valorquestao]
As ligações para um centro de manutenção ocorrem a uma taxa (média) de 2,7 chamadas por
minuto. Qual a probabilidade de que:
%
\begin{enumerate}[(a)]
    \item não mais de quatro chamadas sejam recebidas em um minuto qualquer;
    %
    \item menos de duas chamadas sejam recebidas em um minuto qualquer;
    %
    \item mais de dez chamadas sejam recebidas em um período de 5 minutos.
\end{enumerate}
%
\end{exercicioBanco}


    \noindent\rule{\linewidth}{1pt}  % linha horizontal fina
    
    \phantomsection
    \addcontentsline{toc}{section}{Exercício 14}
    %%%%%%%%%%%%%%%%%%%%%%%%%%%%%%%%%
%
%       EXERCÍCIO
%
%%%%%%%%%%%%%%%%%%%%%%%%%%%%%%%%%

\ifdefstring{\atividade}{prova}{%
    \renewcommand{\valorquestao}{\ValorQ\ pontos}
}%

\begin{exercicioBanco}[\valorquestao]
Suponha que o tempo (em horas) de falha de ventiladores em um computador pessoal possa ser modelado por uma distribuição exponencial, com média de 3300.
%
\begin{enumerate}[(a)]
    \item Qual a proporção de ventiladores que durará no mínimo 10.000 horas?
    %
    \item Qual a proporção de ventiladores que durará no máximo 7.000 horas?
    %
    \item Calcule \(P(X > 10000 \mid X > 8000)\) e \(P (X > 2000)\). Observe que esses resultados são iguais, consequência de uma importante propriedade da distribuição Exponencial chamada de Propriedade da Falta de Memória.
\end{enumerate}
%
\end{exercicioBanco}


    \noindent\rule{\linewidth}{1pt}  % linha horizontal fina
    
    \phantomsection
    \addcontentsline{toc}{section}{Exercício 15}
    %%%%%%%%%%%%%%%%%%%%%%%%%%%%%%%%%
%
%       EXERCÍCIO
%
%%%%%%%%%%%%%%%%%%%%%%%%%%%%%%%%%

\definevalorquestao

\begin{exercicioBanco}[\valorquestao]
Seja \(X \sim \calnN(5,4)\). Determine
%
\begin{center}
    \begin{tabularx}{\textwidth}{XXX}
        (a) \(P(X \le 6)\). &
        (b) \(P(7 < X < 8)\). &
        (c) \(P(X \le -1)\).
    \end{tabularx}
\end{center}
%

\end{exercicioBanco}


    \noindent\rule{\linewidth}{1pt}  % linha horizontal fina
    
    \phantomsection
    \addcontentsline{toc}{section}{Exercício 16}
    %%%%%%%%%%%%%%%%%%%%%%%%%%%%%%%%%
%
%       EXERCÍCIO
%
%%%%%%%%%%%%%%%%%%%%%%%%%%%%%%%%%

\ifdefstring{\atividade}{prova}{%
    \renewcommand{\valorquestao}{\ValorQ\ pontos}
}%

\begin{exercicioBanco}[\valorquestao]
A resistência à compressão de amostras de cimento pode ser modelada por uma distribuição normal com média de 6000 kg/cm\(^{2}\) e um desvio padrão de 100 kg/cm\(^{2}\).
%
\begin{enumerate}[(a)]
    \item Qual a probabilidade da resistência de uma amostra de cimento ser menor que 6250 kg/cm\(^{2}\)?
    %
    \item Qual a probabilidade da resistência de uma amostra de cimento estar entre 5800 kg/cm\(^{2}\) e 5900 kg/cm\(^{2}\)?
    %
    \item Que resistência é excedida por 95\% das amostras de cimento?
\end{enumerate}
%
\end{exercicioBanco}


    \noindent\rule{\linewidth}{1pt}  % linha horizontal fina
    
    \phantomsection
    \addcontentsline{toc}{section}{Exercício 17}
    %%%%%%%%%%%%%%%%%%%%%%%%%%%%%%%%%
%
%       EXERCÍCIO
%
%%%%%%%%%%%%%%%%%%%%%%%%%%%%%%%%%

\definevalorquestao

\begin{exercicioBanco}[\valorquestao]
Seja \(X \sim \mathcal{N}(\mu,\sigma^{2})\). Determine \(a \in \bbR\) tal que \(P(\mu - a\sigma \le X \le \mu + a\sigma) = 0,99\).
%
\end{exercicioBanco}


    \noindent\rule{\linewidth}{1pt}  % linha horizontal fina
    
\end{document}
%*************** FINAL DO DOCUMENTO ***************
